\section{Related Work}\label{sec:related}
% Comparison table and audit for this section: repo .claude_logs/research/novelty-and-related-work.md and
% research/2026-10-20-structural-novelty-audit.md. Bibliography checked 2026-10-28 (notes/reference_check.md).

\subsection{Multimodal Fusion for Segmentation}
Early RGB-X segmentation fused one auxiliary sensor into an RGB stream, for example by cross-modal feature rectification~\cite{cmx2023} or by exchanging tokens between modality-specific transformers~\cite{tokenfusion2022}. CMNeXt~\cite{cmnext2023} generalized this to an arbitrary number of sensors with a hub that aggregates the auxiliary sensors before injecting them into RGB, and introduced DELIVER. Later work improved efficiency and accuracy with adapter-based early fusion~\cite{stitchfusion2024} and pixel-level token exchange~\cite{geminifusion2024}. CAFuser~\cite{cafuser2025} conditions the fusion on the scene condition inferred from the camera image, and DGFusion~\cite{dgfusion2026} uses depth, with depth supervision, to guide the fusion; both report leading results on DELIVER and MUSES. These methods are designed and evaluated with all sensors present. We re-train two of them with degraded inputs and show that their robustness depends mostly on training, not on their fusion design.

\subsection{Robustness to Missing and Degraded Sensors}
A second line of work trains models that tolerate missing sensors. MAGIC~\cite{magic2024,magicpp2024} treats all modalities equally and, during training, ranks modality features by their similarity to an aggregated feature. AnySeg~\cite{anyseg2024} names the problem unimodal bias, the over-reliance on a dominant, faster-to-learn modality, and reduces it by distilling a multimodal teacher into a student trained with random modality dropout; functional-entropy regularization~\cite{zheng2025entropy} targets the same bias. RMMSS~\cite{rmmss2025} combines prototype distillation with feature selection. A recent benchmark~\cite{mmssbench2025} standardized entire-missing, random-missing, and noisy-modality evaluations, which we adopt, and attributed robustness differences to fusion design. Our study differs from these works in its question and its controls: rather than proposing another robust architecture, we take the accuracy-leading fusion models, hold their architecture fixed, change only the training data, and measure how their dependence on each sensor changes across conditions. FuseAdapt~\cite{fuseadapt2026} adapts a frozen backbone with parameter-efficient fusion modules for missing modalities. Our training scheme uses the same building blocks as several of these methods, namely synthetic degradation and distillation from a clean teacher. We do not claim it as a new method; we use it to show where over-reliance comes from and to obtain a model that is both accurate and robust. Methods in this line usually skip a missing input, whereas the benchmark protocol replaces it with zeros; we keep the two families apart when comparing numbers.

\subsection{Adapting Frozen Foundation Models}
Vision foundation models are increasingly used as frozen encoders for dense prediction. The Segment Anything Model~\cite{sam2023} has been adapted to multimodal inputs with adapters that inject auxiliary-sensor features into the RGB stream~\cite{mmsamadapter2025} or with mixtures of LoRA experts~\cite{mlesam}, and MemorySAM~\cite{memorysam2025} reuses the memory mechanism of SAM~2 to fuse modalities. DINOv3~\cite{dinov3_2025} provides dense features and a segmentation adapter derived from ViT-Adapter~\cite{vitadapter2023}, which we reuse for our high-resolution path. Our backbone, LoRA adaptation~\cite{lora2022}, and query head~\cite{mask2former2022} follow this practice; we treat them as the platform for our study, not as contributions.
